%%%%%%%%%%%%%%%%%%%%%%%%%%%%%%%%%%%%
%%%%%%%%%%%%%%%%%%%%%%%%%%%%%%%%%%%%
\section{Academic Integrity}
%%%%%%%%%%%%%%%%%%%%%%%%%%%%%%%%%%%%
%%%%%%%%%%%%%%%%%%%%%%%%%%%%%%%%%%%%
\label{sec:cheating}

(See: {\footnotesize\url{https://www.cmu.edu/policies/student-and-student-life/academic-integrity.html}})

Students in this course are expected to meet CMU's academic integrity standards in full ---
see the full university policy linked above. 
In short, this means submitting your own work, giving credit for others' ideas and contributions, 
and maintaining the trust that makes an academic community possible. 
Students who cannot meet these standards should voluntarily withdraw from this course.


%%%%%%%%%%%%%%%%%%%%%%%%%%%%%%%%%%%%
\subsection{Use of Generative AI in this course}
%%%%%%%%%%%%%%%%%%%%%%%%%%%%%%%%%%%%

To best support your own learning, 
you should complete all graded assignments in this course yourself, 
without any use of generative artificial intelligence (AI). 
Please refrain from using AI tools to generate any content 
(text, video, audio, images, code, etc.) 
for an assignment or classroom exercise. 
Passing off any AI generated content as your own 
(e.g., cutting and pasting content into written assignments, or paraphrasing AI content) 
constitutes a violation of CMU’s academic integrity policy%
\footnote{\url{https://www.cmu.edu/policies/student-and-student-life/academic-integrity.html}}. 
If you have any questions about using generative AI in this course please email or talk to us.

%%%%%%%%%%%%%%%%%%%%%%%%%%%%%%%%%%%%
\subsection{Cheating and Plagiarism}
%%%%%%%%%%%%%%%%%%%%%%%%%%%%%%%%%%%%
In order to deter and detect plagiarism, online tools and other resources may be used in this class.
It is the ethical responsibility of each student to identify the conceptual sources of any work submitted; 
failure to do so is dishonest and is the basis for a charge of cheating or plagiarism, 
which is subject to disciplinary action.

\textbf{Cheating} includes but is not necessarily limited to:
\begin{enumerate}
\item Plagiarism, explained below.
\item Submission of work that is not the student's own for papers, assignments, or exams.
\item Submission or use of falsified data.
\item Theft of or unauthorized access to an exam.
\item Use of an alternate, stand-in or proxy during an examination.
\item Use of unauthorized material including textbooks, notes or computer programs in the preparation of an assignment or during an examination.
\item Supplying or communicating in any way unauthorized information to another student for the preparation of an assignment or during an examination.
\item  Unauthorized collaboration in the preparation of an assignment (see ``Collaboration vs. Cheating,'' Section~\ref{sec:collaboration}, below).
\item Submission of the same work for credit in two courses without obtaining the permission of the instructors beforehand.
\end{enumerate}

\textbf{Plagiarism} includes, but is not limited to, failure to indicate the source with quotation marks or footnotes where appropriate if any of the following are reproduced in the work submitted by a student:
\begin{enumerate}
\item A phrase, written or musical.
\item A graphic element.
\item A proof.
\item Specific language.
\item An idea derived from the work, published or unpublished, of another person.
\end{enumerate}
Any disciplinary actions regarding charges of cheating or plagiarism will follow the procedures of the home university of the student involved.
 
\paragraph{Collaboration vs. Cheating}
\label{sec:collaboration}
Collaboration is defined by Merriam-Webster’s Collegiate Dictionary (10th edition) as 
“to work jointly with others or together, especially in an intellectual endeavor.” 
You are encouraged to discuss the course material, concepts, and homework assignments with other students in the class. 
\textbf{However, each student must eventually submit their own unique work (e.g., homework solutions, final project). }
If any collaboration was used to complete an assignment, 
record the names of the collaborators and the nature of the collaboration. 
Any attempt to submit work that is not the student’s own work will be considered an act of cheating. 
In addition, any student who knowingly supplies their completed homework assignment 
to another student is violating the cheating policy and will be subject to disciplinary action.
 
 \begin{tcolorbox}[colback=red!5,colframe=red!75!black]
ANY VIOLATION OF THIS POLICY WILL NOT BE TOLERATED AND THE PENALTY WILL BE FAILURE IN THE COURSE.
\end{tcolorbox}
 
\textbf{\textit{If you have any questions regarding this policy, contact the instructors.}}

